\documentclass[a4paper,10pt]{article}
\usepackage[utf8]{inputenc}
\usepackage[T1]{fontenc}
\usepackage[french]{babel}
\usepackage[margin=1cm]{geometry} % Marges réduites pour tout faire tenir
\usepackage{xcolor}
\usepackage{tcolorbox}
\usepackage{enumitem}
\usepackage{lmodern}
\usepackage{multicol} % Pour faire des colonnes si besoin

% --- DESIGN & COULEURS ---

% Police sans serif (plus propre pour une fiche)
\renewcommand{\familydefault}{\sfdefault}

% Palette de couleurs
\definecolor{mainBlue}{RGB}{0, 70, 100}      % Bleu foncé (Titres principaux)
\definecolor{softTeal}{RGB}{0, 140, 160}     % Cyan (Sous-titres / Définitions)
\definecolor{alertRed}{RGB}{180, 40, 40}     % Rouge (Attention / Important)
\definecolor{bgGray}{RGB}{245, 247, 250}     % Gris très clair (Fond des boîtes)

% Configuration des boites (tcolorbox)
\tcbset{
    colback=bgGray,           % Couleur de fond
    colframe=mainBlue,        % Couleur du cadre
    coltitle=white,           % Couleur du titre
    fonttitle=\bfseries\large,% Police du titre
    arc=3mm,                  % Arrondi des coins
    boxrule=0.6mm,            % Épaisseur du trait
    left=2mm, right=2mm, top=2mm, bottom=2mm, % Marges internes
    after skip=0.4cm          % Espace après la boite
}

% Commandes rapides pour le texte
\newcommand{\notion}[1]{\textbf{\textcolor{mainBlue}{#1}}}
\newcommand{\important}[1]{\textit{\textcolor{alertRed}{#1}}}
\newcommand{\definition}[1]{\textbf{\textcolor{softTeal}{#1}}}

% --- DÉBUT DU DOCUMENT ---
\begin{document}

% EN-TÊTE
\begin{center}
    \colorbox{mainBlue}{\parbox{\dimexpr\textwidth-2\fboxsep}{
        \centering
        \vspace{0.2cm}
        {\Huge \textbf{\textcolor{white}{FICHE DE RÉVISION : CEJM Thème 1}}} \\
        \vspace{0.1cm}
        {\Large \textcolor{white}{Intégration de l'entreprise dans son environnement}}
        \vspace{0.2cm}
    }}
\end{center}
\vspace{0.3cm}

% ---------------------------------------------------------
% PARTIE 1 : ÉCONOMIE
% ---------------------------------------------------------
\begin{tcolorbox}[title=1. ÉCONOMIE : Agents et Marchés]
    \begin{multicols}{2}
        
    \subsection*{A. Les Agents Économiques}
    Regroupés selon leur fonction principale :
    \begin{itemize}[leftmargin=*]
        \item \notion{Ménages} : Consommation.
        \item \notion{Entreprises} : Production de biens et services \textbf{marchands} (But lucratif).
        \item \notion{État} : Production de services \textbf{non marchands} (financés par impôts).
        \item \notion{Banques} : Financement.
    \end{itemize}

    \subsection*{B. Le Marché}
    Lieu de rencontre Offre/Demande.
    \begin{itemize}[leftmargin=*]
        \item \notion{Loi de l'offre et la demande} :
        \begin{itemize}
            \item Offre > Demande $\rightarrow$ Prix $\searrow$
            \item Demande > Offre $\rightarrow$ Prix $\nearrow$
        \end{itemize}
        \item \notion{Prix d'équilibre} : Offre = Demande.
    \end{itemize}

    \columnbreak

    \subsection*{C. Dysfonctionnements}
    Le marché n'est pas toujours parfait (CPP) :
    \begin{itemize}[leftmargin=*]
        \item \definition{Monopole} : 1 seul offreur.
        \item \definition{Oligopole} : Quelques offreurs (risque d'entente).
        \item \definition{Asymétrie d'information} : Déséquilibre d'info entre vendeur et acheteur (ex: occasion).
        \item \definition{Externalités} :
        \begin{itemize}
            \item \important{Négative} (Pollution) : L'État taxe.
            \item \notion{Positive} (Recherche) : L'État subventionne.
        \end{itemize}
    \end{itemize}

    \end{multicols}
\end{tcolorbox}

% ---------------------------------------------------------
% PARTIE 2 : DROIT
% ---------------------------------------------------------
\begin{tcolorbox}[title=2. DROIT : Le Contrat, colframe=softTeal]
    
    \subsection*{A. Principes Fondateurs}
    \begin{itemize}[leftmargin=*]
        \item \notion{Liberté contractuelle} : Libre de contracter ou non, de choisir son partenaire.
        \item \notion{Force obligatoire} : Le contrat est la loi des parties.
        \item \notion{Bonne foi} : Loyauté obligatoire.
    \end{itemize}

    \vspace{0.2cm}
    \hrule
    \vspace{0.2cm}

    \subsection*{B. Validité du contrat (Art. 1128 C. Civ)}
    3 conditions cumulatives pour éviter la \important{Nullité} :
    \begin{enumerate}[leftmargin=*]
        \item \definition{Consentement} : Doit être libre. Sans vices :
        \begin{itemize}
            \item \textbf{Erreur} : Involontaire.
            \item \textbf{Dol} : Tromperie volontaire.
            \item \textbf{Violence} : Pression physique/morale.
        \end{itemize}
        \item \definition{Capacité} : Être majeur capable.
        \item \definition{Contenu} : Licite (légal) et Certain (existant).
    \end{enumerate}

    \subsection*{C. Clauses de sécurisation}
    \begin{itemize}
        \item \notion{Réserve de propriété} : Transfert de propriété seulement après paiement complet.
        \item \notion{Clause pénale} : Fixe l'indemnité à l'avance en cas de problème.
        \item \notion{Clause résolutoire} : Annulation automatique si inexécution.
    \end{itemize}
\end{tcolorbox}

% ---------------------------------------------------------
% PARTIE 3 : MANAGEMENT
% ---------------------------------------------------------
\begin{tcolorbox}[title=3. MANAGEMENT : Performance et RSE, colframe=alertRed]
    \begin{multicols}{2}
    
    \subsection*{A. Finalités (P. Drucker)}
    \begin{itemize}[leftmargin=*]
        \item \notion{Économique} : Profit, pérennité.
        \item \notion{Sociale} : Salariés (bien-être).
        \item \notion{Sociétale} : RSE, Environnement.
    \end{itemize}

    \subsection*{B. Parties Prenantes (Freeman)}
    Tout acteur affecté par l'entreprise :
    \begin{itemize}[leftmargin=*]
        \item \definition{Internes} : Salariés, Actionnaires.
        \item \definition{Externes} : Clients, État, ONG.
    \end{itemize}

    \columnbreak 

    \subsection*{C. La Performance}
    \begin{itemize}[leftmargin=*]
        \item \notion{Efficacité} : Atteindre l'objectif.
        \item \notion{Efficience} : Optimiser les moyens.
        \item \notion{Globale} : Éco + Sociale + Sociétale.
    \end{itemize}
    \textit{Outil :} Tableau de bord prospectif (Kaplan).
    \end{multicols}
\end{tcolorbox}

% ---------------------------------------------------------
% LEXIQUE
% ---------------------------------------------------------
\begin{tcolorbox}[title=LEXIQUE EXPRESS, colback=white, colframe=black]
    \footnotesize
    \begin{description}[style=multiline, leftmargin=3cm]
        \item[\textcolor{mainBlue}{Agent Économique}] Centre de décision économique indépendant.
        \item[\textcolor{mainBlue}{Contrat}] Accord de volontés créant des obligations (donner, faire, ne pas faire).
        \item[\textcolor{mainBlue}{RSE}] Responsabilité Sociétale des Entreprises (impact social/écolo).
        \item[\textcolor{mainBlue}{Externalité}] Effet de l'activité d'un agent sur un tiers sans compensation financière.
        \item[\textcolor{mainBlue}{Nullité}] Anéantissement rétroactif du contrat (il n'a jamais existé).
    \end{description}
\end{tcolorbox}

\end{document}